\documentclass[a4paper,11pt]{article}
\usepackage[utf8]{inputenc}
\usepackage[T1]{fontenc}
\usepackage[french]{babel}
\usepackage[left=1cm,right=1cm,top=.5cm,bottom=0cm]{geometry}
\usepackage{enumitem}
\usepackage{hyperref}
\usepackage{xcolor}
\usepackage{titlesec}
\usepackage{fontawesome5}
\usepackage{lmodern}
\usepackage[normalem]{ulem}

% --- Couleurs ---
\definecolor{primary}{RGB}{35, 55, 123}
\definecolor{secondary}{RGB}{80, 80, 80}

% --- Configuration des liens ---
\hypersetup{colorlinks=true, linkcolor=primary, filecolor=primary, urlcolor=primary}

% --- Formatage des titres ---
\titleformat{\section}{\large\bfseries\color{primary}\uppercase}{}{0em}{}[\titlerule]
\titlespacing{\section}{0pt}{8pt}{5pt}

% --- Commandes personnalisées ---
\newcommand{\resumeItem}[1]{\item\small{#1}}
\newcommand{\resumeSubheading}[4]{
  \vspace{-1pt}\item
    \begin{tabular*}{0.97\textwidth}[t]{l@{\extracolsep{\fill}}r}
      \textbf{#1} & \textit{\small #2} \\
      \textit{\small #3} & \textit{\small #4} \\
    \end{tabular*}\vspace{-5pt}
}

\begin{document}

% --- EN-TÊTE ---
\begin{center}
    {\Large \textbf{Loïc Aron MBASSI EWOLO}} \\ \vspace{4pt}
    \textbf{\large Stage Ingénieur IA - Analyse de Défauts} \\
    \textit{\small Disponible dès \textbf{Avril 2026} pour 4 à 6 mois} \\ \vspace{4pt}
    \small
    \faMapMarker\ Cergy, France (Mobile) \quad
    \faPhone\ (+33) 7 59 52 33 98 \quad
    \faEnvelope\ \href{mailto:loicaron.mbassiewolo@ensea.fr}{loicaron.mbassiewolo@ensea.fr} \\ \vspace{2pt}
    \faLinkedin\ \href{https://www.linkedin.com/in/loic-aron-mbassi-ewolo-3942a9246/}{linkedin} \quad
    \faGithub\ \href{https://github.com/Nameless0l}{github} \quad
    \faGlobe\ \href{https://mbassiloic.tech}{mbassiloic.tech}
\end{center}

\vspace{-6pt}

% --- PROFIL ---
\noindent\small{Élève-ingénieur en double diplôme \textbf{ENSEA/Polytechnique Yaoundé}, spécialisé en traitement du signal et apprentissage automatique. Expérience en détection d'anomalies sur signaux (ECG) et en classification d'images. Je recherche un stage au \textbf{CEA Gramat} pour appliquer mes compétences en Python et Machine Learning à l'analyse de défauts.}

% --- FORMATION ---
\section{Formation}
\begin{itemize}[leftmargin=*, label=\tiny$\bullet$, nosep]
    \item \textbf{ENSEA (École Nationale Supérieure de l'Électronique et de ses Applications)} \hfill \textit{Cergy, France | 2025 -- Présent} \\
    \small{Ingénieur 2A, Double Diplôme - \textbf{Traitement du Signal} et Intelligence Artificielle.} \\
    \textit{\small{Cours : Traitement numérique du signal, Filtrage, Statistiques, Apprentissage automatique.}}
    \vspace{2pt}
    \item \textbf{École Polytechnique de Yaoundé (1\textsuperscript{ère} école d'ingénieur CEMAC)} \hfill \textit{Cameroun | 2021 -- 2025} \\
    \small{Diplôme d'Ingénieur de Conception (Master 2) : \textbf{Mathématiques Appliquées}, Algorithmique, Génie Logiciel.}
    \item \textbf{Université de Yaoundé I} \hfill \textit{Cameroun | 2020 -- 2022} \\
    \small{Licence Informatique \& Mathématiques : Analyse, Algèbre Linéaire, Probabilités.}
\end{itemize}

% --- COMPÉTENCES TECHNIQUES ---
\section{Compétences Techniques}
\begin{itemize}[leftmargin=*, nosep]
    \item \small{\textbf{IA \& Data Science :} \textbf{Python} (NumPy, Pandas, Scikit-learn), PyTorch, TensorFlow/Keras, Computer Vision.}
    \item \small{\textbf{Traitement du Signal :} Filtrage numérique, Analyse spectrale, Séries temporelles, Extraction de features.}
    \item \small{\textbf{Mathématiques :} Algèbre Linéaire (SVD), Probabilités/Statistiques, Analyse Numérique, Optimisation.}
    \item \small{\textbf{Outils :} Git, Linux, Docker, Jupyter, Matplotlib/Seaborn.}
\end{itemize}

% --- EXPÉRIENCES RECHERCHE ---
\section{Expériences en Recherche \& IA}
\begin{itemize}[leftmargin=*, label={-}]

    \resumeSubheading
      {Assistant de Recherche - Généralisation des Réseaux de Neurones}{Juin 2025 -- Sept. 2025}
      {MILA (Institut Québécois d'IA) - Remote}{Québec, Canada}
      \begin{itemize}[leftmargin=0.4cm, nosep, label=\tiny$\bullet$]
        \item \small{Travaux sur le phénomène de « Grokking » (généralisation tardive) dans les réseaux de neurones.}
        \item \small{Implémentation de pipelines de tests sous \textbf{PyTorch}, reproduction d'expériences SOTA.}
        \item \small{Analyse mathématique du comportement des modèles face au surapprentissage.}
      \end{itemize}
    \vspace{2pt}

    \resumeSubheading
      {Projet de Recherche UROP - Analyse de Signaux ECG}{2024}
      {École Polytechnique de Yaoundé (Mentors Google DeepMind)}{Cameroun}
      \begin{itemize}[leftmargin=0.4cm, nosep, label=\tiny$\bullet$]
        \item \small{Développement d'un modèle de \textbf{détection d'anomalies} cardiaques à partir d'images d'ECG.}
        \item \small{Prétraitement : filtrage du bruit, extraction de features sur séries temporelles.}
        \item \small{Classification avec TensorFlow/Keras, évaluation sur données réelles.}
      \end{itemize}

\end{itemize}

% --- PROJETS IA ---
\section{Projets en Machine Learning}
\begin{itemize}[leftmargin=*, label={-}]

\item \textbf{Violence Detection - Classification Vidéo} \hfill \textit{Compétition | 2022}
\vspace{-2pt}
\begin{itemize}[leftmargin=0.4cm, nosep, label=\tiny$\bullet$]
    \item \small{Modèle de \textbf{détection d'anomalies} (violence) dans des flux vidéo en temps réel.}
    \item \small{Technologies : TensorFlow, OpenCV. \textbf{2\textsuperscript{ème} place} de la compétition.}
\end{itemize}
\vspace{3pt}

\item \textbf{Women Safe - Détection de Contenus Sexistes (NLP)} \hfill \textit{Compétition MLPC | 2023}
\vspace{-2pt}
\begin{itemize}[leftmargin=0.4cm, nosep, label=\tiny$\bullet$]
    \item \small{État de l'art et comparaison de modèles (BERT, GPT-2) pour la détection de misogynie.}
    \item \small{Étude des biais éthiques dans les modèles de langage.}
\end{itemize}
\vspace{3pt}

\item \textbf{Urban Waste Management - Prédiction \& Optimisation} \hfill \textit{Hackathon CITS | 2024}
\vspace{-2pt}
\begin{itemize}[leftmargin=0.4cm, nosep, label=\tiny$\bullet$]
    \item \small{Modèle de prédiction des besoins de collecte de déchets à partir de données IoT.}
    \item \small{Algorithmes d'optimisation de routes. \textbf{1\textsuperscript{er} prix} contre 75 équipes.}
\end{itemize}
\vspace{3pt}

\item \textbf{ClassSync - Reconnaissance Faciale Temps Réel} \hfill \textit{Projet Académique | 2024}
\vspace{-2pt}
\begin{itemize}[leftmargin=0.4cm, nosep, label=\tiny$\bullet$]
    \item \small{Système de détection et reconnaissance faciale pour appel automatique.}
    \item \small{Technologies : OpenCV, Face Recognition. Présenté au salon national GETEC.}
\end{itemize}

\end{itemize}

% --- CERTIFICATIONS & LANGUES ---
\section{Certifications, Leadership \& Langues}
\begin{itemize}[leftmargin=*, nosep]
    \item \textbf{Certifications :} Supervised ML (DeepLearning.AI/Stanford), Python for Data Science (IBM).
    \item \textbf{Leadership :} Chef de la Cellule IA (Polytechnique Yaoundé), animation de workshops ML.
    \item \textbf{Hackathons :} \textbf{1\textsuperscript{ère} Place} IndabaX Africa (IA Santé), \textbf{1\textsuperscript{ère} Place} CITS National.
    \item \textbf{Langues :} Français (Natif), Anglais (Technique/Scientifique).
    \item \textbf{Note Bac Physique :} 19.25/20.
\end{itemize}

\end{document}
